% f(x)=x1+.2x2，x=(.8,0)，L=sqrt(1.04)，认证r=.7。
\begin{tikzpicture}[figure picture,foundation picture]
 \path[use as bounding box] (0,0) rectangle (110,80);
 \begin{scope}[xshift=10mm]
 \fill[FigureBlueFill] (29,7)--(74,7)--(74,67)--(17,67)--cycle;
 \draw[draw=FigureStroke,line width=1.1pt] (29,7)--(17,67);
 \draw[draw=LancetGreen,fill=LancetGreen!12!white,line width=1pt] (43,37) circle (17.5);
 % 坐标轴置于区域与认证球的填充之上，保持整条轴可追踪。
 \draw[foundation axis] (4,37)--(78,37);
 \node[anchor=north] at (78,33) {$x_1$};
 \draw[foundation axis] (23,5)--(23,70) node[above] {$x_2$};
 \fill[FigureInk] (43,37) circle (1.4);
 \node[anchor=south west] at (44,39) {$\bm x=(0.8,0)$};
 \draw[draw=FigureMuted,line width=.6pt] (43,37)--(43,54.5);
 \node[anchor=west,inner sep=1.5pt] at (44,49) {$r=0.7$};
 \draw[draw=FigureBlue,fill=white,line width=1pt] (66,62) circle (1.5);
 \node[anchor=south] at (66,65) {$f>0$};
 \draw[draw=FigureRed,line width=1pt] (9,45.5)--(12,48.5) (9,48.5)--(12,45.5);
 \node[anchor=south] at (10.5,52.5) {$f<0$};
 \node[anchor=north] at (29,5) {$f=0$};
 \end{scope}
\end{tikzpicture}%
